\documentclass[../cours_1re.tex]{subfiles}
\begin{document}
\setcounter{chapter}{5}
\chapter{Modèles discrets et suites affines}
\label{ch:06}
\begin{objectif}
Modéliser un processus par une récurrence affine ; expliciter son évolution ; étudier un équilibre et un cumul. Les suites affines et les modèles financiers sont des approfondissements.
\end{objectif}
\section{Le sens du modèle}
Une population fictive perd $20\,\%$ de ses individus à chaque période, puis reçoit $30$ nouveaux individus. Le modèle est $u_{n+1}=0{,}8u_n+30$.
Si l'arrivée précède la perte, il devient $u_{n+1}=0{,}8(u_n+30)$. L'ordre des opérations appartient à la modélisation.
Les valeurs non entières peuvent représenter un effectif moyen théorique ; elles ne sont pas des fractions d'individus observés.
\section{Centrer autour d'un équilibre}
\begin{theoreme}
Pour $u_{n+1}=au_n+b$ avec $a\ne1$, poser $\ell=b/(1-a)$. Alors
\[
u_n=\ell+a^n(u_0-\ell).
\]
Si $a=1$, la suite est arithmétique : $u_n=u_0+nb$.
\end{theoreme}
\begin{demonstration}
L'égalité $\ell=a\ell+b$ caractérise l'équilibre. En la soustrayant à la récurrence,
$u_{n+1}-\ell=a(u_n-\ell)$. La suite $(u_n-\ell)$ est géométrique ; le chapitre 5 donne la formule. Le cas $a=1$ se lit directement dans la relation initiale.
\end{demonstration}
\begin{exemple}
Pour $u_0=50$ et $u_{n+1}=0{,}8u_n+30$, $\ell=150$ et $u_n=150-100(0{,}8)^n$.
La suite croît, reste inférieure à $150$ et s'en rapproche. Son écart à l'équilibre est exactement $100(0{,}8)^n$.
\end{exemple}
\section{Stabilité et oscillations}
Si $|a|<1$, les écarts successifs sont multipliés en valeur absolue par un facteur inférieur à $1$ ; ils tendent vers zéro, comme démontré au chapitre 7.
Pour $-1<a<0$, ils alternent de signe : convergence n'implique pas monotonie.
Pour $|a|>1$, un écart initial non nul s'amplifie. Si $a=-1$, l'écart alterne sans diminuer ; sauf au point fixe, il n'y a pas de limite.

Deux suites de même modèle et de conditions initiales $u_0,v_0$ vérifient
$u_n-v_n=a^n(u_0-v_0)$. Cela mesure exactement la sensibilité à l'état initial.
\section{Une récurrence affine sans point fixe}
Une autre méthode donne, pour $n\ge1$,
\[
u_n=a^nu_0+b\sum_{k=0}^{n-1}a^k.
\]
\begin{demonstration}
Au rang $1$, la formule donne $au_0+b$. Si elle est vraie au rang $n$, multiplier par $a$ et ajouter $b$ donne
$a^{n+1}u_0+b\sum_{k=1}^{n}a^k+b$, soit la formule au rang $n+1$.
\end{demonstration}
Cette écriture reste correcte pour $a=1$. Chaque apport $b$ est pondéré par le nombre de périodes pendant lesquelles il a été multiplié par $a$.
\section{Annuités constantes : un exemple mathématique}
On modélise une dette initiale $D_0>0$, un taux périodique $i\ge0$ et un remboursement $A$ en fin de période, après application des intérêts :
$D_{n+1}=(1+i)D_n-A$.
Pour $i>0$,
\[
D_n=(1+i)^nD_0-A\,\frac{(1+i)^n-1}{i}.
\]
Imposer $D_N=0$ donne
\[
A=\frac{iD_0}{1-(1+i)^{-N}}.
\]
\begin{demonstration}
Appliquer la formule précédente avec $a=1+i$, $b=-A$, puis isoler $A$ dans $D_N=0$. Pour $i=0$, revenir à $D_n=D_0-nA$, d'où $A=D_0/N$ ; la formule avec division par $i$ ne s'applique pas directement.
\end{demonstration}
\begin{exemple}
Avec un taux nul, $D_0=1200$ et $N=12$, $A=100$. Avec un taux positif, une partie de chaque versement rémunère les intérêts : pour diminuer la dette dès le premier versement, il faut $A>iD_0$.
Le modèle ignore frais, arrondis contractuels et variations de taux ; sa portée est algébrique.
\end{exemple}
\section{Tour de Hanoï et croissance exponentielle}
\approfondissement Pour déplacer $n$ disques empilés entre trois tiges, sans jamais poser un grand disque sur un petit, noter $H_n$ le nombre minimal de mouvements.
Avant de déplacer le plus grand disque, il faut libérer les $n-1$ autres ; après, il faut les replacer. Cela demande au moins $2H_{n-1}+1$ mouvements. La stratégie récursive réalise exactement ce nombre, donc
$H_n=2H_{n-1}+1$, $H_0=0$.
\begin{demonstration}
$H_n+1=2(H_{n-1}+1)$, donc $H_n+1=2^n$ et $H_n=2^n-1$. Il s'agit bien d'un minimum : la borne inférieure et la stratégie réalisant cette borne ont toutes deux été établies.
\end{demonstration}
Le passage de $n$ à $n+1$ double presque le travail. Cette croissance explique pourquoi un algorithme correct peut devenir impraticable.

\begin{finchapitre}
\section{Utilité et prolongements}
Soustraire un équilibre transforme une récurrence affine en multiplication. C'est une première technique de changement de variable, réutilisable lorsqu'une formule semble compliquée.

En licence, une récurrence vectorielle s'écrit $U_{n+1}=AU_n+B$ ; le comportement des puissances de la matrice $A$ dépend de ses valeurs propres. En master, la stabilité de systèmes dynamiques et de schémas numériques généralise la condition $|a|<1$. Le calcul exact de l'écart constitue déjà une estimation d'erreur.
\begin{retenir}
Commencer par préciser l'ordre des événements. Chercher un équilibre puis étudier l'écart ; ne pas supposer qu'un équilibre est automatiquement attractif.
\end{retenir}
\end{finchapitre}
\end{document}
